%%=============================================================================
%% Samenvatting
%%=============================================================================

% TODO: De "abstract" of samenvatting is een kernachtige (~ 1 blz. voor een
% thesis) synthese van het document.
%
% Een goede abstract biedt een kernachtig antwoord op volgende vragen:
%
% 1. Waarover gaat de bachelorproef?
% 2. Waarom heb je er over geschreven?
% 3. Hoe heb je het onderzoek uitgevoerd?
% 4. Wat waren de resultaten? Wat blijkt uit je onderzoek?
% 5. Wat betekenen je resultaten? Wat is de relevantie voor het werkveld?
%
% Daarom bestaat een abstract uit volgende componenten:
%
% - inleiding + kaderen thema
% - probleemstelling
% - (centrale) onderzoeksvraag
% - onderzoeksdoelstelling
% - methodologie
% - resultaten (beperk tot de belangrijkste, relevant voor de onderzoeksvraag)
% - conclusies, aanbevelingen, beperkingen
%
% LET OP! Een samenvatting is GEEN voorwoord!

%%---------- Nederlandse samenvatting -----------------------------------------
%
% TODO: Als je je bachelorproef in het Engels schrijft, moet je eerst een
% Nederlandse samenvatting invoegen. Haal daarvoor onderstaande code uit
% commentaar.
% Wie zijn bachelorproef in het Nederlands schrijft, kan dit negeren, de inhoud
% wordt niet in het document ingevoegd.

\IfLanguageName{english}{%
\selectlanguage{dutch}
\chapter*{Samenvatting}
\lipsum[1-4]
\selectlanguage{english}
}{}

%%---------- Samenvatting -----------------------------------------------------
% De samenvatting in de hoofdtaal van het document

\chapter*{\IfLanguageName{dutch}{Samenvatting}{Abstract}}

Freshfrozen For Pets is een West-Vlaamse KMO die verse maaltijden voor honden en katten produceert. De productieleveringen komen binnen op/in leeggoed die na gebruik teruggestuurd moeten worden naar de leverancier. Het beheer van dat leeggoed gebeurt vandaag volledig in Microsoft Excel, met per leverancier een eigen spreadsheet. Hoewel Excel laagdrempelig is, brengt deze werkwijze een aantal terugkerende problemen met zich mee: invoer is niet afdwingbaar, waarden worden destructief overschreven, er is geen historiek of audit-spoor en het ERP-systeem van MAPA (de leverancier van het systeem), waarin de productieproducten wel geregistreerd worden, kent het leeggoed niet. Daardoor ontbreekt zowel een overzicht per leverancier als de mogelijkheid om foute boekingen achteraf te traceren.

Vanuit deze probleemstelling is de onderzoeksvraag ontstaan: Hoe kan het beheer van leeggoed bij Freshfrozen For Pets efficiënter georganiseerd worden, met als focus om tijdbesparing, foutenreductie, procesmatige gebruiksvriendelijkheid en overzichtelijkheid te verbeteren? Drie fases zijn uitgeschreven om een antwoord op deze vraag te vormen. Fase 1 onderzoekt een directe uitbreiding van het bestaande, Progress-gebaseerde ERP-systeem. Fase 2 onderzoekt een middleware-opstelling waarbij een afzonderlijk leeggoedsysteem via een point-to-point-integratie met het ERP-systeem communiceert. Fase 3 betreft een volledige standalone-applicatie, uitgewerkt als een functionele proof-of-concept met Blazor WebAssembly, ASP.NET en PostgreSQL.

Het onderzoek start met een intern onderzoek bij Freshfrozen. Er werden interviews gehouden met de CEO, de productiechef en een financieel bediende, een enquête afgenomen bij de dagelijkse gebruikers en er was overleg met MAPA. Daaruit bleek dat een concrete uitvoering van fase 1 en fase 2 in de praktijk niet mogelijk was omwille van licentie- en codebase-restricties op het Progress-platform, waardoor deze twee fases theoretisch werden uitgewerkt. De PoC van fase 3 werd wel volledig geïmplementeerd en iteratief gevalideerd met een eindgebruiker.

Fase 1 scoort het sterkst op tijdsbesparing en foutenreductie omdat de leeggoedingave op de bestaande receptiebon afdwingbaar wordt en het ERP zelf de referentiecode aanlevert. Ook is dit de omgeving waarmee de eindgebruikers al het meest vertrouwd zijn. De schermen die uitgewerkt zijn in fase 3 zijn wel effectief wat Freshfrozen nodig heeft. De optimale oplossing combineert beide werelden: de modulestructuur en interfacekeuzes uit fase 3 verder laten uitwerken binnen het ERP-systeem van MAPA. Op die manier worden de afdwingbaarheid en de afwezigheid van datasilo's uit fase 1 verenigd met de overzichtelijkheid uit fase 3.

Voor Freshfrozen biedt dit onderzoek een concreet vervolgtraject in samenwerking met MAPA. 
Een goede aanbeveling is dat elke case anders is, dus deze drie methodes van implementatie kunnen als leidraad gebruikt worden maar zijn geen definitief antwoord op elk probleem. 
